% GENERATED FILE. Edit canonical lessons, then run npm run generate:books.
% canonical-source-hash: fnv1a64:e7dd49a17642cd8c
% canonical-lessons: SA-C01-namaste, SA-C01-namaskara, SA-S02-letter-na, SA-C01-dhanyavada, SA-C01-svagatam, SA-S200-letter-da, SA-C01-am, SA-C01-na, SA-C01-practice, SA-S05-sign-virama

\chapter{Greetings}
\label{ch:greetings}
\hlchaptermodality{1}

\begin{chapteropening}
\textbf{By the end of this chapter:} \emph{I can greet someone in Sanskrit, recognize the reply, and choose an appropriate response.}

Complete SA-C01-practice, the canonical closing checkpoint, and demonstrate the chapter promise: greet someone in Sanskrit, recognize the reply, and choose an appropriate response.
\end{chapteropening}

\section[namaste]{\sk{नमस्ते} (namaste) — \textquotedblleft{}I bow to you\textquotedblright{}}

\label{lesson:SA-C01-namaste}



\begin{quote}
The first Sanskrit word — and the original that Hindi, Marathi, Punjabi, and Bengali all wore down into their own greetings.
\end{quote}

\begin{sounds}
Four even beats: \emph{na · ma · s · te}. The \emph{s} closes its own beat and leans straight into \emph{te} with no vowel between them. Nothing here is strange to an English mouth — say it at conversational speed and it is yours.
\end{sounds}

\begin{cousinweb}
\emph{namaste} = \emph{namas} (\textquotedblleft{}a bow, obeisance\textquotedblright{}) + \emph{te} (\textquotedblleft{}to you\textquotedblright{}) — literally \textquotedblleft{}\textbf{a bow to you}.\textquotedblright{} This is the \emph{source}: the word every Indo-Aryan \emph{namaste} descends from. \emph{Namas} is from the root √nam \textquotedblleft{}to bend, bow\textquotedblright{} (PIE *\emph{nem-}). And \emph{te} \textquotedblleft{}to you\textquotedblright{} is a family heirloom — the same Indo-European pronoun behind Latin \emph{tē}, Greek \emph{se}, and archaic English \textbf{thee}. Say \emph{namas-te} and you speak a word Cicero and Chaucer would each half-recognise.
\end{cousinweb}

\begin{culture}
Palms pressed together (the gesture is \emph{añjali}), a slight bow: hands and word say the same thing. \emph{Namaste} greets the person before you as worthy of a bow — hello, and on parting, goodbye. It is the greeting Sanskrit gave to half of South Asia.
\end{culture}

\subsection*{Your turn}
Say these aloud:

\begin{itemize}
\raggedright
  \item read it — na · ma · s · te
  \item the two parts (\emph{namas} \textquotedblleft{}a bow\textquotedblright{} + \emph{te} \textquotedblleft{}to you\textquotedblright{})
  \item the English cousin of \emph{te} (\textbf{thee})
\end{itemize}

\begin{tcolorbox}[breakable,colback=teal!4,colframe=teal!35!black,title={Before you move on}]
What does \emph{namaste} literally mean, and what English pronoun is \emph{te} cousin to? (\textquotedblleft{}A bow to you\textquotedblright{}; \textbf{thee}.)
\end{tcolorbox}
\section[namaskāraḥ]{\sk{नमस्कारः} (namaskāraḥ) — the fuller form}

\label{lesson:SA-C01-namaskara}



\begin{quote}
Same bow, different grammar — and your first case-ending.
\end{quote}

\begin{sounds}
\emph{na · mas · kā · raḥ}. The \emph{ā} is the long vowel, held about twice as long as a short \emph{a}. The final \emph{-ḥ} is a soft breath after the vowel, almost an echo of it: the word ends on a whispered \emph{ha}, not a stop.
\end{sounds}

\begin{cousinweb}
\emph{namaskāraḥ} = \emph{namas} (\textquotedblleft{}a bow\textquotedblright{}) + \emph{kāra} (\textquotedblleft{}the making, the doing,\textquotedblright{} from the root √kṛ \textquotedblleft{}to do/make\textquotedblright{}). So \emph{namaskāra} is \textquotedblleft{}the \textbf{making} of a bow\textquotedblright{} — an act, where \emph{namaste} is an address. That root √kṛ \textquotedblleft{}to do\textquotedblright{} is among Sanskrit's most productive (it gives \emph{karma}, \textquotedblleft{}a thing done\textquotedblright{}); it is PIE \emph{k\textsuperscript{w}er-}. The final \textbf{visarga} \emph{-aḥ} marks the masculine singular — your first case-ending.
\end{cousinweb}

\begin{grammarlens}[title={Grammar Lens}]
\emph{Namaste} and \emph{namaskāra} share the same first half, \emph{namas}, but end differently because they do different grammatical jobs — one takes a pronoun (\emph{te} \textquotedblleft{}to you\textquotedblright{}), the other builds a noun (\emph{-kāra} \textquotedblleft{}the making of\textquotedblright{}). Sanskrit words are transparently assembled from roots and endings; once you see the seams, long words stop being frightening.
\end{grammarlens}

\subsection*{Your turn}
Say these aloud:

\begin{itemize}
\raggedright
  \item read it — na · mas · kā · raḥ
  \item how it differs in sense from \emph{namaste} (\textquotedblleft{}the making of a bow\textquotedblright{} vs. \textquotedblleft{}a bow to you\textquotedblright{})
  \item what the final visarga \emph{-aḥ} marks (masculine singular)
\end{itemize}

\begin{tcolorbox}[breakable,colback=teal!4,colframe=teal!35!black,title={Before you move on}]
\emph{Namaskāra} is built from which two parts, and what does its visarga ending tell you? (\emph{Namas} \textquotedblleft{}a bow\textquotedblright{} + \emph{kāra} \textquotedblleft{}the making of\textquotedblright{}; masculine singular.)
\end{tcolorbox}
\section[na]{\sk{न} — one character, met inside words you already say}

\label{lesson:SA-S02-letter-na}



\begin{quote}
A horizontal head-line (shirorekhā) runs across the top; letters hang beneath it like laundry on a line.

One character this time. One only — and you have been saying it for pages without knowing which mark on the page it was.
\end{quote}

\subsection*{Script you'll notice: \sk{न}}
\textbf{\sk{न}} — \emph{na}.

It is a \textbf{consonant}, and in this script a consonant is never bare: it comes with an \emph{a} already in it. So it is not \emph{n}, it is \textbf{na}.

What it is made of:

\begin{itemize}
\raggedright
  \item a clockwise left loop joined to a rightward shoulder
  \item a top-to-bottom right stem
  \item the top shirorekhā
\end{itemize}

You already say these words, and \sk{न} is one of the shapes inside them — the rest of their shapes are still ahead of you, so here the character stands on its own:

\begin{itemize}
\raggedright
  \item \emph{namaste} — hello (lit. \textquotedblleft{}a bow to you\textquotedblright{})
  \item \emph{namaskāraḥ} — hello (lit. \textquotedblleft{}the making of a bow\textquotedblright{})
\end{itemize}

\subsection*{Writing: \sk{न}}
\begin{itemize}
\raggedright
  \item \textbf{1.} start at the loop's inner-right curve, sweep down and clockwise around the small opening, then continue right along the shoulder without lifting
  \item \textbf{2.} lift and draw the right stem top-to-bottom
  \item \textbf{3.} lift and draw the top shirorekhā left-to-right
\end{itemize}

\textbf{Pen lifts: 2.} The pen comes up 2 times and no more.

\begin{quote}
verified three-stroke teaching form; everyday handwriting may join or simplify the loop differently
\end{quote}

\begin{quote}
This is one attested teaching order and not a national standard — handwriting here is taught with school-to-school variation. Source: Opiaterein, ‘Deva-\sk{न}-order.gif’, strokes 1–3, Wikimedia Commons, 11 May 2009.
\end{quote}

\subsection*{Your turn}
\begin{itemize}
\raggedright
  \item \emph{Look:} at this, and find \sk{न} in it
\end{itemize}

\begin{quote}
\sk{न}
\end{quote}

\begin{itemize}
\raggedright
  \item \emph{Trace it:} \sk{न} three times, saying \emph{na} as you finish each one
  \item \emph{Look:} back at any page of this chapter and find \sk{न} once more
\end{itemize}

\begin{tcolorbox}[breakable,colback=teal!4,colframe=teal!35!black,title={Before you move on}]
Which character is this — \sk{न}? What sound does it carry? (\emph{\textbf{na}}.) Name one word you already say that contains it.
\end{tcolorbox}
\section[dhanyavādaḥ]{\sk{धन्यवादः} (dhanyavādaḥ) — \textquotedblleft{}thank you\textquotedblright{}}

\label{lesson:SA-C01-dhanyavada}



\begin{quote}
The full, un-worn Sanskrit shape of a \textquotedblleft{}thank you\textquotedblright{} you already met to the east.
\end{quote}

\begin{sounds}
\emph{dha · nya · vā · daḥ}. The opening \emph{dh} is one sound, not two — a \emph{d} released with a puff of breath behind it. \emph{nya} is likewise a single run. It closes on the same soft breathed \emph{-ḥ} as \emph{namaskāraḥ}.
\end{sounds}

\begin{cousinweb}
\emph{dhanyavādaḥ} = \emph{dhanya} (\textquotedblleft{}worthy, deserving,\textquotedblright{} from \emph{dhana} \textquotedblleft{}wealth\textquotedblright{}) + \emph{vāda} (\textquotedblleft{}a speaking, a statement,\textquotedblright{} from the root √vad \textquotedblleft{}to speak\textquotedblright{}). Literally \textquotedblleft{}a saying of 'you are worthy.'\textquotedblright{} This is the \textbf{source} of the thanks you met to the east — Hindi \emph{dhanyavād}, Marathi and Punjabi likewise, Bengali \emph{dhônyobad} — here in its full Sanskrit shape, visarga and all.
\end{cousinweb}

\begin{culture}
Classical Sanskrit thanks are more often shown in deed and honorific than said with a set formula; \emph{dhanyavāda} is the plain, dignified \textquotedblleft{}thank you,\textquotedblright{} and its survival across every modern northern language shows how durable a good word can be. Note you already knew all its pieces — only the ending (\emph{-aḥ}) is new.
\end{culture}

\subsection*{Your turn}
Say these aloud:

\begin{itemize}
\raggedright
  \item read it — dha · nya · vā · daḥ
  \item its two roots (\emph{dhanya} \textquotedblleft{}worthy\textquotedblright{} + \emph{vāda} \textquotedblleft{}a saying\textquotedblright{})
  \item a modern descendant (Hindi \emph{dhanyavād} / Bengali \emph{dhônyobad})
\end{itemize}

\begin{tcolorbox}[breakable,colback=teal!4,colframe=teal!35!black,title={Before you move on}]
\emph{Dhanyavāda} is built from which two parts, and what modern greeting descends from it? (\emph{Dhanya} \textquotedblleft{}worthy\textquotedblright{} + \emph{vāda} \textquotedblleft{}a saying\textquotedblright{}; Hindi \emph{dhanyavād}, etc.)
\end{tcolorbox}
\section[svāgatam]{\sk{स्वागतम्} (svāgatam) — \textquotedblleft{}welcome\textquotedblright{}}

\label{lesson:SA-C01-svagatam}



\begin{quote}
One word that is \emph{cognate} with English \textquotedblleft{}welcome\textquotedblright{} — not just a translation of it.
\end{quote}

\begin{sounds}
\emph{svā · ga · tam}. It opens on \emph{sv-} said as one push, carries the long \emph{ā} in that first beat, and ends on a bare \emph{m} with no vowel after it — the mouth simply closes.
\end{sounds}

\begin{cousinweb}
\emph{svāgatam} = \emph{su} (\textquotedblleft{}well, good\textquotedblright{}) + \emph{āgata} (\textquotedblleft{}come,\textquotedblright{} from \emph{ā} \textquotedblleft{}toward\textquotedblright{} + the root √gam \textquotedblleft{}to go/come\textquotedblright{}) — \textquotedblleft{}\textbf{well come},\textquotedblright{} which is exactly what English \emph{welcome} means too. Two deep Indo-European threads surface:

\begin{itemize}
\raggedright
  \item \emph{su-} \textquotedblleft{}good\textquotedblright{} is the twin of \textbf{Greek \emph{eu-}} — the \emph{eu} of \emph{euphoria}, \emph{eulogy}, \emph{euphemism}. Sanskrit \emph{su-} and Greek \emph{eu-} are the same prefix.
  \item the root √gam \textquotedblleft{}to go/come\textquotedblright{} is PIE \emph{g\textsuperscript{w}em-} — the very root of English \textbf{come} (and of Latin \emph{venīre}, hence \emph{venue}, \emph{convene}).
\end{itemize}

So \emph{svāgatam} and \emph{welcome} are not just parallel ideas but partly \textbf{cognate} ones.
\end{cousinweb}

\begin{culture}
\emph{Su} + \emph{āgata} fuses into \emph{svāgata} by \textbf{sandhi} — Sanskrit's rule that \emph{u} before a vowel becomes the glide \emph{v}. You will see sandhi everywhere; it is why Sanskrit words seem to melt into one another. The neuter ending \emph{-am} makes it \textquotedblleft{}a welcome / it is welcomed.\textquotedblright{}
\end{culture}

\subsection*{Your turn}
Say these aloud:

\begin{itemize}
\raggedright
  \item read it — svā · ga · tam
  \item the two parts (\emph{su} \textquotedblleft{}well\textquotedblright{} + \emph{āgata} \textquotedblleft{}come\textquotedblright{})
  \item the English word that is \emph{su-}'s Greek twin (\emph{eu-}, as in \emph{euphoria})
  \item the English cousin of the root √gam (\textbf{come})
\end{itemize}

\begin{tcolorbox}[breakable,colback=teal!4,colframe=teal!35!black,title={Before you move on}]
\emph{Svāgatam} literally means what, and name two of its far-flung cousins. (\textquotedblleft{}Well come\textquotedblright{} = welcome; Greek \emph{eu-} (from \emph{su-}) and English \emph{come} (from √gam).)
\end{tcolorbox}
\section[da]{\sk{द} — one character, met inside words you already say}

\label{lesson:SA-S200-letter-da}



\begin{quote}
Before the new one, one you have already met: \sk{न}. Say what it does.

One character this time. One only — and you have been saying it for pages without knowing which mark on the page it was.
\end{quote}

\subsection*{Script you'll notice: \sk{द}}
\textbf{\sk{द}} — \emph{da}.

It is a \textbf{consonant}, and in this script a consonant is never bare: it comes with an \emph{a} already in it. So it is not \emph{d}, it is \textbf{da}.

What it is made of:

\begin{itemize}
\raggedright
  \item a short top-to-bottom stem
  \item an outer body joined to an inward curl and down-right tail
  \item the top shirorekhā
\end{itemize}

You already say this word, and \sk{द} is one of the shapes inside it — the rest of its shapes are still ahead of you, so here the character stands on its own:

\begin{itemize}
\raggedright
  \item \emph{dhanyavādaḥ} — thank you
\end{itemize}

\subsection*{Writing: \sk{द}}
\begin{itemize}
\raggedright
  \item \textbf{1.} draw the short stem top-to-bottom
  \item \textbf{2.} lift at the stem's lower junction, sweep left through the upper shoulder and around the outer body, curl inward and clockwise through the loop, then continue down-right through the tail without lifting
  \item \textbf{3.} lift and draw the top shirorekhā left-to-right
\end{itemize}

\textbf{Pen lifts: 2.} The pen comes up 2 times and no more.

\begin{quote}
verified three-stroke teaching form; another learner source stages the outer body and curl-tail separately
\end{quote}

\begin{quote}
This is one attested teaching order and not a national standard — handwriting here is taught with school-to-school variation. Source: Opiaterein, ‘Deva-\sk{द}-order.gif’, strokes 1–3, Wikimedia Commons, 8 May 2009.
\end{quote}

\subsection*{Your turn}
\begin{itemize}
\raggedright
  \item \emph{Look:} at this, and find \sk{द} in it
\end{itemize}

\begin{quote}
\sk{द}
\end{quote}

\begin{itemize}
\raggedright
  \item \emph{Trace it:} \sk{द} three times, saying \emph{da} as you finish each one
  \item \emph{Look:} back at any page of this chapter and find \sk{द} once more
\end{itemize}

\begin{tcolorbox}[breakable,colback=teal!4,colframe=teal!35!black,title={Before you move on}]
Which character is this — \sk{द}? What sound does it carry? (\emph{\textbf{da}}.) Name one word you already say that contains it.
\end{tcolorbox}
\section[ām]{\sk{आम्} (ām) — \textquotedblleft{}yes\textquotedblright{}}

\label{lesson:SA-C01-am}



\begin{quote}
You can greet, thank, and welcome. Now the smallest reply of all.
\end{quote}

\begin{sounds}
\emph{ām} — one beat. A long \emph{ā}, then the lips close on \emph{m} with no vowel after it. Say it and stop.
\end{sounds}

\begin{culture}
Classical Sanskrit, like Latin, leans on more than a bare \textquotedblleft{}yes\textquotedblright{}: one answers by echoing the verb, or with \emph{ām}, \emph{evam} (\textquotedblleft{}so it is\textquotedblright{}), or \emph{bāḍham} (\textquotedblleft{}certainly\textquotedblright{}). The blunt yes/no habit of English is not universal — these old languages preferred to \emph{say what they meant}. \emph{Ām} is the plain agreement, and the one to reach for first.
\end{culture}

\subsection*{Your turn}
Say these aloud:

\begin{itemize}
\raggedright
  \item yes — \emph{ām}
  \item two other ways to agree (\emph{evam} \textquotedblleft{}so it is\textquotedblright{}, \emph{bāḍham} \textquotedblleft{}certainly\textquotedblright{})
  \item what Sanskrit prefers to a bare yes (echoing the verb)
\end{itemize}

\begin{tcolorbox}[breakable,colback=teal!4,colframe=teal!35!black,title={Before you move on}]
How do you say yes, and what does Sanskrit often do instead of saying it? (\emph{Ām}; echo the verb, or answer with \emph{evam} or \emph{bāḍham}.)
\end{tcolorbox}
\section[na]{\sk{न} (na) — \textquotedblleft{}no, not\textquotedblright{}}

\label{lesson:SA-C01-na}



\begin{quote}
You have \emph{ām}. Here is its opposite — and it may be the oldest thing in this whole book.
\end{quote}

\begin{sounds}
\emph{na} — one short beat, the \emph{a} closer to English \emph{cup} than to \emph{car}.
\end{sounds}

\begin{cousinweb}
\emph{Na} (\textquotedblleft{}no, not\textquotedblright{}) you have already met wearing other clothes: it is Proto-Indo-European *\emph{ne} \textquotedblleft{}not,\textquotedblright{} \textbf{the same negative} as the Latin \emph{nōn} of the last track, as English \textbf{no}, \textbf{not}, \textbf{none}, as German \emph{nein}, and as Hindi \emph{nahīṇ}. Of every word in this book, \emph{na} may reach furthest unchanged — one syllable spoken, in some form, from Iceland to India for five thousand years.
\end{cousinweb}

\subsection*{Your turn}
Say these aloud:

\begin{itemize}
\raggedright
  \item yes and no — \emph{ām} / \emph{na}
  \item the English cousins of \emph{na} (no, not, none)
  \item the Latin cousin you already know (\emph{nōn})
\end{itemize}

\begin{tcolorbox}[breakable,colback=teal!4,colframe=teal!35!black,title={Before you move on}]
What is the ancient root of \emph{na}, and name three of its living cousins. (PIE *\emph{ne} \textquotedblleft{}not\textquotedblright{}; English \emph{no/not/none}, Latin \emph{nōn}, German \emph{nein}, Hindi \emph{nahīṇ} — any three.)
\end{tcolorbox}
\section[Practice]{Chapter 1 — Practice and Recall}

\label{lesson:SA-C01-practice}



\begin{quote}
No new words — just the five, and the roots that reach both east and west.
\end{quote}

\subsection*{Your turn}
Say these aloud:

\begin{itemize}
\raggedright
  \item hello, \textquotedblleft{}a bow to you\textquotedblright{} — \emph{namaste}
  \item hello, \textquotedblleft{}the making of a bow\textquotedblright{} — \emph{namaskāraḥ}
  \item thank you — \emph{dhanyavādaḥ}
  \item welcome, \textquotedblleft{}well come\textquotedblright{} — \emph{svāgatam}
  \item yes and no — \emph{ām} / \emph{na}
\end{itemize}

\subsection*{How to answer: the taproot both ways}
Sanskrit points east \emph{and} west. Trace each:

Say these aloud:

\begin{itemize}
\raggedright
  \item which Sanskrit word do Hindi/Marathi/Punjabi/Bengali's greetings come from? (\emph{namas} / \emph{namaste})
  \item which Greek prefix is \emph{su-} \textquotedblleft{}good\textquotedblright{} the twin of? (\emph{eu-})
  \item which English verb is the root √gam cousin to? (\emph{come})
  \item which Latin word (last track) is \emph{na} cousin to? (\emph{nōn})
\end{itemize}

\begin{grammarlens}[title={What you've built: the roots you now carry}]
\begin{itemize}
\raggedright
  \item \emph{namas} \textquotedblleft{}a bow\textquotedblright{} $\to$ every Indo-Aryan hello; root √nam \textquotedblleft{}to bend\textquotedblright{}
  \item \emph{te} \textquotedblleft{}to you\textquotedblright{} $\to$ English \textbf{thee}, Latin \emph{tē}
  \item \emph{su-} \textquotedblleft{}good\textquotedblright{} $\to$ Greek \textbf{eu-} (euphoria, eulogy)
  \item √gam \textquotedblleft{}come\textquotedblright{} $\to$ English \textbf{come}, Latin \emph{venīre}
  \item \emph{na} \textquotedblleft{}not\textquotedblright{} $\to$ English \textbf{no/not/none}, Latin \emph{nōn}, German \emph{nein}
\end{itemize}
\end{grammarlens}

\begin{tcolorbox}[breakable,colback=teal!4,colframe=teal!35!black,title={Before you move on}]
In 1786 Sir William Jones said Sanskrit, Latin, and Greek shared \textquotedblleft{}a stronger affinity than could have been produced by accident.\textquotedblright{} Name two words from this chapter that show it. (Any two: \emph{te}$\leftrightarrow$\emph{thee}/\emph{tē}, \emph{su-}$\leftrightarrow$\emph{eu-}, √gam$\leftrightarrow$\emph{come}, \emph{na}$\leftrightarrow$\emph{nōn}.)
\end{tcolorbox}
\section[virāma]{\sk{◌्} — one character, met inside words you already say}

\label{lesson:SA-S05-sign-virama}



\begin{quote}
Before the new one, the ones you have already met: \sk{न} · \sk{द}. Say what each of them does.

One character this time. One only — and you have been saying it for pages without knowing which mark on the page it was.
\end{quote}

\subsection*{Script you'll notice: \sk{◌्}}
\textbf{\sk{◌्}} — \emph{no sound of its own — it takes one away}.

It is a \textbf{vowel-killer}. Every consonant in this script arrives with an \emph{a} already inside it — that is what an abugida is. This mark takes the \emph{a} back, leaving the bare consonant. It is how the script writes two consonants in a row.

Where it sits: a stroke below that removes the inherent 'a' and forms a conjunct with the next consonant.

Worked through: \textbf{\sk{न}} + \textbf{\sk{◌्}} = \textbf{\sk{न्}} — \emph{n}.

You already say these words, and \sk{◌्} is one of the shapes inside them — the rest of their shapes are still ahead of you, so here the character stands on its own:

\begin{itemize}
\raggedright
  \item \emph{namaste} — hello (lit. \textquotedblleft{}a bow to you\textquotedblright{})
  \item \emph{namaskāraḥ} — hello (lit. \textquotedblleft{}the making of a bow\textquotedblright{})
  \item \emph{dhanyavādaḥ} — thank you
  \item \emph{svāgatam} — welcome (lit. \textquotedblleft{}well come\textquotedblright{})
\end{itemize}

\subsection*{Writing: \sk{◌्} — copy what you see}
Put your pen on \sk{◌्} and follow its line. Copy the shape you can see — slowly, and larger than it is printed.

\begin{quote}
This book does not yet tell you \textbf{where to start the character or which way to travel}. That is a real thing, taught with real variation from school to school, and it is not written down here until it can be written down with a source. Copying what is in front of you needs no such source, and it is how the shape gets into your hand in the meantime.
\end{quote}

\subsection*{Your turn}
\begin{itemize}
\raggedright
  \item \emph{Look:} at this, and find \sk{◌्} in it
\end{itemize}

\begin{quote}
\sk{◌्}
\end{quote}

\begin{itemize}
\raggedright
  \item \emph{Trace it:} \sk{◌्} three times, and each time say what it does: \textbf{it kills the built-in \emph{a}}
  \item \emph{Look:} back at any page of this chapter and find \sk{◌्} once more
\end{itemize}

\begin{tcolorbox}[breakable,colback=teal!4,colframe=teal!35!black,title={Before you move on}]
Which character is this — \sk{◌्}? What does it do to the consonant it sits on? (\textbf{Takes away the built-in \emph{a}.}) Name one word you already say that contains it.
\end{tcolorbox}
